% =====================================================================
%  IEEE-STYLE CONFERENCE REPORT TEMPLATE  (self-contained)
%  MLE-AI-201 Machine Learning -- Final Project
%
%  This template MIMICS the IEEE two-column conference look using only
%  standard LaTeX packages. It does NOT depend on IEEEtran.cls, so it
%  compiles anywhere (tectonic, pdflatex, Overleaf) with nothing extra.
%
%  Build:   tectonic report.tex      (recommended)
%     or:   pdflatex report.tex ; pdflatex report.tex
% =====================================================================
\documentclass[10pt,twocolumn,a4paper]{article}

% ---- Geometry: IEEE-like margins / two narrow columns ----------------
\usepackage[a4paper,top=0.75in,bottom=1.0in,left=0.68in,right=0.68in,columnsep=0.25in]{geometry}

% ---- Times-like font (standard, ships with every TeX distro) ---------
\usepackage{mathptmx}          % Times text + math
\usepackage[T1]{fontenc}
\usepackage[utf8]{inputenc}    % allows accented / Azerbaijani chars

% ---- Common, dependency-free packages --------------------------------
\usepackage{graphicx}
\usepackage{amsmath,amssymb}
\usepackage{booktabs}
\usepackage{array}
\usepackage{titlesec}
\usepackage{caption}
\usepackage{cite}
\usepackage[hidelinks]{hyperref}
\usepackage{url}

% ---- IEEE-style section headings -------------------------------------
\renewcommand{\thesection}{\Roman{section}}
\renewcommand{\thesubsection}{\Alph{subsection}}
\titleformat{\section}[block]{\centering\normalsize\scshape}{\thesection.}{0.5em}{}
\titleformat{\subsection}[block]{\normalsize\itshape}{\thesubsection.}{0.5em}{}
\titlespacing*{\section}{0pt}{1.4ex plus 1ex minus .2ex}{1ex}
\titlespacing*{\subsection}{0pt}{1.0ex plus 1ex minus .2ex}{0.6ex}

% ---- Abstract / Index-Terms in the IEEE boldface-dash style ----------
\newcommand{\abstractkw}[1]{\noindent\textbf{\textit{#1}}}

\captionsetup{font=small,labelsep=period}
\setlength{\parindent}{1em}

% ---- Project helpers ---------------------------------------------------
\usepackage{xcolor}
\usepackage{enumitem}
\setlist{nosep,leftmargin=1.4em}
\graphicspath{{figures/}}

% Anything still open is printed in red, so nothing unfinished can be missed.
% Before submission:  grep -c '\\todo{' report.tex   must print 0.
\newcommand{\todo}[1]{\textcolor{red}{\textbf{[TODO: #1]}}}

% Tables and single numbers are written by `python -m src.run_all` into
% tables/. Run it before building this file. Nothing in them is typed by hand.
\input{tables/numbers}
% A generated table, shrunk only if it is wider than the column.
\newcommand{\gentable}[1]{\resizebox{\ifdim\width>\linewidth\linewidth\else\width\fi}{!}{\input{tables/#1}}}
% A generated figure, or a visible note if the model behind it does not exist yet.
\newcommand{\genfigure}[2][\linewidth]{%
  \IfFileExists{figures/#2.png}{\includegraphics[width=#1]{#2.png}}{%
    \fbox{\parbox[c][2.4cm][c]{0.9\linewidth}{\centering\small\todo{figure
    \texttt{\detokenize{#2}} is produced by the pipeline once the model exists}}}}}

% =====================================================================
\begin{document}

% ------------------------- TITLE BLOCK --------------------------------
\twocolumn[
  \begin{@twocolumnfalse}
  \begin{center}
    {\LARGE\bfseries A Decision Tree and a Pegasos SVM from Scratch:\\
     Predicting Baku Property Prices from bina.az Listings\par}
    \vspace{1.2em}
    % Authors: Surname, Given name --- alphabetical by surname.
    {\large
      Alimammadova, Nigar\quad Dilavarli, Farid\quad Gojayev, Nihad\quad
      Karimova, Ulkar\par}
    \vspace{0.5em}
    {\itshape MLE-AI-201 Machine Learning, AI Academy --- Cohort II\par}
    \vspace{1.4em}
  \end{center}
  \begin{quote}
    \small
    \abstractkw{Abstract---}%
    We predict property prices in Baku from \NumListings{} cleaned sale
    listings scraped from bina.az, as a regression task (price) and as a
    derived classification task (premium versus standard, split at the
    median training price). The two core models, a decision tree and a
    linear soft-margin support-vector machine trained with the Pegasos
    sub-gradient method, are implemented from scratch in NumPy and
    compared with scikit-learn under matched settings, the same splits and
    the same seeds. The cleaning removes two target-leakage columns and a
    less obvious leak: 36\% of the raw rows are repeat scrapes of a
    listing that is already in the file. Our tree reproduces
    scikit-learn's to the third decimal (test $R^2$ 0.781 on the log
    price, F1 0.912 for the tier) and gives the same tier call on 99.9\%
    of held-out listings; our SVM reaches F1 0.849 against 0.852 for
    \texttt{SGDClassifier}, with 95.7\% identical calls. The tree beats
    the linear SVM by six F1 points, because price depends on area and
    location jointly. An error analysis shows where the models fail:
    the most expensive 1\% of listings, land plots, commercial property
    and districts with few training listings. One command reproduces every
    number, table and figure:
    \href{https://github.com/nigar001/ML_project}{\texttt{github.com/nigar001/ML\_project}}
    (tag \texttt{v1.0-final}).

    \vspace{0.6em}
    \abstractkw{Index Terms---}decision tree, support-vector machine,
    Pegasos, regression, classification, real-estate price prediction.
  \end{quote}
  \vspace{1.0em}
  \end{@twocolumnfalse}
]

% =====================================================================
\section{Introduction}
Asking prices on bina.az, the largest property platform in Azerbaijan,
depend on size, location, building type and condition in ways that are
neither linear nor uniform across the market. We use a public scrape of
its sale listings~\cite{dataset} to study two questions with models we
wrote ourselves: how well can the price of a listing be predicted
(Task~A, regression), and can a listing be placed in the upper or lower
half of the market (Task~B, classification)?

The aim is to understand two standard algorithms from the inside and to
compare them honestly with a mature library, not to maximise one score.
Our contributions are:
\begin{itemize}
  \item a documented cleaning pipeline for a raw scraped file, including
        the removal of target leakage and of repeat scrapes;
  \item a decision tree (Gini, entropy and variance reduction) and a
        Pegasos SVM with a random-Fourier-feature extension, both in
        NumPy only;
  \item metrics, cross-validation and a benchmark harness written from
        scratch, with scikit-learn~\cite{sklearn} used only as the
        baseline and to check our numbers;
  \item an error analysis by price band, property category and district
        coverage;
  \item a seeded pipeline in which one command regenerates every table
        and figure of this report.
\end{itemize}

\section{Background}
A decision tree~\cite{cart} partitions the feature space with
axis-aligned threshold splits chosen greedily to reduce an impurity
measure, and predicts the majority class or the mean target of the
training points in each leaf. It needs no feature scaling and captures
interactions, but an unconstrained tree memorises its training data, so
depth and leaf size must be limited.

A soft-margin support-vector machine~\cite{svm} finds a separating
hyperplane with a wide margin while tolerating violations. Its primal
form is an unconstrained minimisation of a regularised hinge loss, which
Pegasos~\cite{pegasos} solves by stochastic sub-gradient descent with a
step size of $1/(\lambda t)$, with no quadratic-programming solver. A
linear SVM can be made non-linear by mapping the inputs to random
Fourier features~\cite{rff}, whose inner products approximate a Gaussian
kernel.

\section{Data and Preprocessing}
\subsection{The raw file}
The dataset~\cite{dataset} (file dated 3~February 2025, downloaded
12~October 2026) has 100{,}775 rows and 51 columns. All prices are in
AZN; 99.7\% of the rows are in Baku. It mixes seven property categories:
new-build and old-build flats, houses, land, commercial property, offices
and garages. Raw prices run from 11~AZN to 600~million AZN and are
strongly right-skewed (skewness 252), while their logarithm is close to
symmetric (skewness 0.16), see Fig.~\ref{fig:price}.

\begin{figure}[t]
  \centering
  \genfigure{price_distribution}
  \caption{Distribution of the cleaned price and of its base-10 logarithm.}
  \label{fig:price}
\end{figure}

\subsection{Leakage}
Two columns are arithmetic functions of the target and are dropped:
\texttt{total\_price} equals \texttt{price} in 99.99\% of the rows, and
\texttt{unit\_price} multiplied by the area reproduces the price to within
0.3\%. Features are taken from an explicit list, so identifiers, owner and
agency names, the address text and the columns recorded after publication
(\texttt{views}, \texttt{vip}, \texttt{featured}) never reach a model.

A second leak is not a column. The site was scraped repeatedly, and the
file holds only 64{,}454 distinct listings: 36\% of the rows are another
scrape of a listing that is already present, some listings up to 17
times. No two rows are fully identical, because each scrape has its own
id, so dropping exact duplicates removes nothing. A random split of rows
would place copies of one flat in both the training and the test data
(48\% of the test rows, in our check), and the test score would partly
measure memory. We keep one row per listing URL, the most recent scrape.

\subsection{Cleaning decisions}
\begin{itemize}
  \item \textit{Floor.} The text ``7 / 9'' is split into the floor and
        the number of floors in the building.
  \item \textit{Area.} Land is listed in sot and buildings in m$^2$; sot
        is converted (1~sot = 100~m$^2$).
  \item \textit{Outliers.} We keep prices from 1{,}000 to 10~million AZN
        and areas from 15 to 3{,}000~m$^2$ (land up to 50{,}000~m$^2$).
        These limits remove 206 of the 64{,}454 distinct listings: 8
        priced below 1{,}000~AZN (values such as 11 or 250~AZN cannot be
        the sale price of a property), 63 above 10~million AZN, 36 under
        15~m$^2$, too small for a dwelling, and 99 above the area cap.
        Land gets a higher cap because plots of several hectares are
        genuine listings.
  \item \textit{Missing values.} Rows without a price or an area are
        dropped. Other missing numbers are set to 0 and missing categories
        to ``Unknown''. These are constants, so no statistic of the data,
        and in particular none of the test split, enters the features.
  \item \textit{Encoding.} Seven numeric features (area, rooms, floor,
        floors in the building, land area, latitude, longitude) and six
        categorical ones (category, repair, bill of sale, mortgage,
        district, city) are used; the categorical ones are one-hot encoded.
\end{itemize}

\subsection{Tasks and split}
Task~A predicts the logarithm of the price. Task~B predicts a binary tier:
premium if the price is above the median of the \emph{training} prices
(\TierThreshold~AZN), standard otherwise. The listings are split once,
with a fixed seed, into training, validation and test parts of
70/15/15\%, stratified so that each part has the same share of premium
listings. Table~\ref{tab:data} gives the sizes.

\begin{table}[t]
  \centering
  \caption{Dataset after cleaning.}
  \label{tab:data}
  \gentable{data_summary}
\end{table}

\section{Method}
\subsection{Decision tree (from scratch)}
For a node with class proportions $p_k$ the tree uses the Gini impurity
or the entropy, and for regression the variance of the target:
\begin{equation}
  G = 1-\sum_k p_k^2,\quad
  H = -\sum_k p_k\log_2 p_k,\quad
  I_{\mathrm{reg}} = \tfrac{1}{N}\sum_i (y_i-\bar{y})^2 .
\end{equation}
A split of a parent node into children $c$ with $N_c$ points is scored by
its impurity decrease,
\begin{equation}
  \Delta I = I(\mathrm{parent}) - \sum_{c}\tfrac{N_c}{N}\,I(c),
\end{equation}
and the tree takes the feature and threshold with the largest decrease.
Candidate thresholds are the midpoints between consecutive distinct sorted
values of a feature. Growth is recursive and stops at \texttt{max\_depth},
below \texttt{min\_samples\_split} points, when a child would have fewer
than \texttt{min\_samples\_leaf} points, or when the best decrease is
below \texttt{min\_impurity\_decrease}. A leaf predicts the majority class
(with class frequencies as probabilities) or the mean target. The
importance of a feature is the total impurity decrease of the splits that
use it, weighted by the share of training points reaching them and
normalised to sum to one.

The split search sorts each feature once per node and obtains the impurity
of every candidate threshold from cumulative class counts (classification)
or cumulative sums of $y$ and $y^2$ (regression), so a node costs
$O(n\log n)$ per feature instead of $O(n^2)$; a test checks this fast
search against a plain reference that scans every threshold. Features are
visited in a seeded random order and the first best threshold is kept, so
ties between equally good splits are broken reproducibly, though not in
scikit-learn's order. A midpoint that rounds onto the upper of two
neighbouring values falls back to the lower one, so no child is ever
empty, and, as in scikit-learn, \texttt{min\_impurity\_decrease} is
compared with the decrease weighted by the node's share of all training
rows. Unit tests check the impurities against hand calculations, that the
tree finds the same splits as scikit-learn on a small fully-determined
example, that each of the four stopping settings changes the tree, and that
\texttt{predict\_proba} returns leaf frequencies.

\subsection{Support-vector machine (from scratch)}
With labels $y_i\in\{-1,+1\}$ the linear soft-margin SVM minimises
\begin{equation}
  J(\mathbf{w},b)=\tfrac{\lambda}{2}\lVert\mathbf{w}\rVert^{2}
  +\tfrac{1}{n}\sum_{i=1}^{n}\max\!\big(0,\,1-y_i(\mathbf{w}^{\top}\mathbf{x}_i+b)\big).
  \label{eq:svm}
\end{equation}
Pegasos~\cite{pegasos} draws a random example $i$ at step $t$, sets
$\eta_t = 1/(\lambda t)$ and updates
\begin{equation}
  \mathbf{w}\leftarrow(1-\eta_t\lambda)\,\mathbf{w}
  +\eta_t\,y_i\mathbf{x}_i\;\mathbf{1}\big[y_i(\mathbf{w}^{\top}\mathbf{x}_i+b)<1\big].
\end{equation}
The margin width is $2/\lVert\mathbf{w}\rVert$, and the support vectors
are the training points with $y_i(\mathbf{w}^{\top}\mathbf{x}_i+b)\le 1$.
The hinge loss is sensitive to scale, so features are standardised with
the mean and standard deviation of the training rows only.
Our implementation folds the bias into $\mathbf{w}$ by appending a
constant feature (so $b$ is regularised too), averages the sub-gradient
over mini-batches of 32 examples drawn with replacement, projects
$\mathbf{w}$ onto the ball $\lVert\mathbf{w}\rVert\le 1/\sqrt{\lambda}$
after each step as in~\cite{pegasos}, and trains for 20 passes over the
training rows. The objective is recorded every 200 steps;
Fig.~\ref{fig:objective} shows it falling and flattening.

\begin{figure}[t]
  \centering
  \genfigure[0.92\linewidth]{objective_ours_svm_pegasos}
  \caption{Training objective $J$ of our SVM (Task B, final fit).}
  \label{fig:objective}
\end{figure}

For the non-linear extension we use random Fourier features~\cite{rff},
$z(\mathbf{x})=\sqrt{2/D}\,\cos(\mathbf{W}\mathbf{x}+\mathbf{u})$ with
Gaussian $\mathbf{W}$ and uniform $\mathbf{u}$, and run the same linear
solver on $z(\mathbf{x})$. Unlike kernelised Pegasos this keeps one weight
vector instead of one coefficient per training point, which matters with
tens of thousands of listings.
The wider SVM study (\texttt{svm\_experiments.py}, run by the same command)
uses $D=500$ and sweeps the bandwidth $\gamma$ together with $\lambda$ on
the validation split. The best random-feature model reached a validation
F1 of 0.809, against 0.825 for the linear SVM in the same study, so the
extension did not pay off on these features and the linear SVM remains our
main model. That study works on a reduced feature set of 54 columns and its
own split, so its numbers are not comparable digit by digit with the
tables below.

\section{Experimental Setup}
\textit{Metrics.} Task~A is scored by RMSE, MAE and $R^2$ on the log
price, and by MAE in AZN after transforming predictions back. Task~B is
scored by precision, recall and F1 for the premium class, by accuracy,
and by ROC-AUC computed from ranks with tied scores sharing their average
rank. The metrics are our own NumPy code; tests check them against
hand-worked examples and against scikit-learn to twelve significant
digits.

\textit{Baselines and matched settings.} Our tree is compared with
\texttt{DecisionTreeRegressor} and \texttt{DecisionTreeClassifier} given
the same criterion and the same four stopping settings. Our SVM is
compared with \texttt{SGDClassifier} with hinge loss and
$\alpha=\lambda$, which minimises the same objective~(\ref{eq:svm}) with
the same kind of solver; its features are standardised inside a
scikit-learn \texttt{Pipeline}. Table~\ref{tab:settings} lists the
settings, which come from \SettingsSource.
\textit{Tuning} (\texttt{src/tuning.py}) uses the training
and validation parts only. For the SVM we sweep
$\lambda\in\{10^{-5},\dots,10^{-1}\}$ with 20 passes each; validation F1
peaks at $\lambda=10^{-3}$ (0.840). For the tree, on a seeded sample of
20{,}000 training rows, we sweep \texttt{max\_depth} from 2 to 15 for Gini
and entropy (Task~B) and for variance (Task~A), then
\texttt{min\_samples\_leaf} $\in\{1,5,20,50\}$ at the chosen depth. Gini
and entropy stay within 0.01 F1 of each other at every depth; Gini peaks
slightly higher and is kept. One set of stopping settings serves both
tasks, so the depth is chosen on Task~A, the tree's main task. The
resulting curves are in Fig.~\ref{fig:tuning}.

\begin{figure*}[t]
  \centering
  \genfigure[0.32\textwidth]{tuning_tree_depth_task_a}\hfill
  \genfigure[0.32\textwidth]{tuning_tree_depth_task_b}\hfill
  \genfigure[0.32\textwidth]{tuning_svm_lambda_f1}
  \caption{Validation curves on the validation split. Left: our tree,
  Task A, $R^2$ against depth. Middle: our tree, Task B, F1 against depth
  for Gini and entropy (the dashed line marks Task B's own best depth).
  Right: our SVM, F1 against $\lambda$.}
  \label{fig:tuning}
\end{figure*}

\begin{table}[t]
  \centering
  \caption{Settings used for every reported number.}
  \label{tab:settings}
  \gentable{settings}
\end{table}

\textit{Protocol.} Nothing from held-out rows is used for a decision.
(i)~The tier threshold is the median of the training prices and is fixed
before any model is fitted. (ii)~Settings are chosen on the validation
part. (iii)~Uncertainty comes from \NumFolds-fold stratified
cross-validation on the training and validation rows; inside each fold the
threshold and the scaling statistics are recomputed from that fold's
training part. (iv)~Each model is then refitted on the training and
validation rows and scored once on the test part. All randomness is seeded
with 42, and two runs of the pipeline give identical numbers. Times are
wall-clock seconds on one laptop (Intel Core i7-13650HX).

\section{Results}
\subsection{Headline results}
Tables~\ref{tab:resultsA} and~\ref{tab:resultsB} give the test-split
metrics of every model, with the cross-validated mean and standard
deviation of the main metric.
Our tree reproduces scikit-learn's to the third decimal on both
tasks: test $R^2$ 0.781 for both on the log price, and F1 0.912 with
AUC 0.969 for both on the tier. Our SVM is within 0.003 F1 of
\texttt{SGDClassifier}. The tree is the stronger model for the tier
(F1 0.912 against 0.849), and its cross-validated scores vary by less
than 0.01 across folds.

\begin{table*}[t]
  \centering
  \caption{Task A, price regression on the log price: test split, and
  cross-validated $R^2$ (mean $\pm$ std over \NumFolds{} folds).}
  \label{tab:resultsA}
  \gentable{results_task_a}
\end{table*}

\begin{table*}[t]
  \centering
  \caption{Task B, price tier: test split, and cross-validated F1
  (mean $\pm$ std over \NumFolds{} folds).}
  \label{tab:resultsB}
  \gentable{results_task_b}
\end{table*}

Fig.~\ref{fig:roc} shows the ROC curves on the test split. The tree is
ahead of the linear SVM over almost the whole range: the tier depends on area and
location jointly, which a tree expresses with a few splits and a single
hyperplane cannot.

\begin{figure}[t]
  \centering
  \genfigure[0.92\linewidth]{roc_curves}
  \caption{ROC curves on the test split, Task B.}
  \label{fig:roc}
\end{figure}

\subsection{Agreement with scikit-learn}
Table~\ref{tab:agreement} compares each of our models with its library
counterpart on one held-out fold. Exact agreement is not expected. Two
trees can differ when two splits have the same impurity decrease, because
each implementation breaks the tie in its own order (scikit-learn visits
features in a seeded random order), and because scikit-learn stores
features as 32-bit floats, which moves thresholds in the last digits. Two
SVMs trained by stochastic steps stop near the optimum of~(\ref{eq:svm}),
not at it, so we compare the objective value they reach.
The trees grow to the same depth with almost the same
number of leaves and share the root split: area $\le 100.25$~m$^2$ for
the price and $\le 92.15$~m$^2$ for the tier, which scikit-learn stores as
92.1500015 because of its 32-bit floats. They give the same tier on
99.87\% of held-out listings, and their price predictions differ by 0.005
on the log scale on average; the few larger gaps come from the handful of
nodes where a tie or a rounded threshold sends a listing to a different
leaf. The SVMs agree on 95.7\% of listings. Ours reaches $J=0.392$
against 0.387: \texttt{SGDClassifier} runs until its own convergence test,
while ours stops after a fixed 20 passes, slightly short of the optimum.

\begin{table*}[t]
  \centering
  \caption{Our models against scikit-learn on one held-out fold.}
  \label{tab:agreement}
  \gentable{agreement}
\end{table*}

\subsection{Margin and support vectors}
Table~\ref{tab:margin} reports the learned $\lVert\mathbf{w}\rVert$, the
margin width and the number of support vectors. The classes overlap
heavily near the median price, so a large share of the training points
lies on or inside the margin.

\begin{table}[t]
  \centering
  \caption{Linear SVMs on the training and validation rows.}
  \label{tab:margin}
  \gentable{svm_margin}
\end{table}

\subsection{What the models use}
Fig.~\ref{fig:importance} shows the impurity-decrease importances of the
tree for Task~A, and Fig.~\ref{fig:weights} the largest weights of the
linear SVM on standardised features. The tree relies mostly on the area
and then on the coordinates; the one-hot district columns add little once
latitude and longitude are available. The SVM cannot split on
coordinates, so besides rooms, area and category it puts weight on
individual districts instead.

\begin{figure}[t]
  \centering
  \genfigure{importances_ours_tree_task_a}
  \caption{Feature importances of our tree, Task A.}
  \label{fig:importance}
\end{figure}

\begin{figure}[t]
  \centering
  \genfigure{svm_weights_ours_svm_pegasos}
  \caption{Largest weights of our linear SVM (standardised features).
  Positive weights push towards the premium tier.}
  \label{fig:weights}
\end{figure}

\subsection{Speed}
Table~\ref{tab:timing} gives fit and predict times at three training
sizes. scikit-learn's tree is compiled code, while ours is NumPy driven
from Python, so a large constant-factor gap is expected.
Our tree fits about ten times more slowly than
scikit-learn's (2.3~s against 0.21~s on 20{,}000 rows for Task~B), but
its time grows almost linearly with the number of rows (0.11, 0.51 and
2.3~s for 1{,}000, 5{,}000 and 20{,}000 rows), as expected from sorting
each feature once per node. Our SVM fits 20{,}000 rows in 0.9~s.

\begin{table*}[t]
  \centering
  \caption{Fit and predict time in seconds by number of training rows.}
  \label{tab:timing}
  \gentable{timing}
\end{table*}

\section{Discussion and Limitations}
\subsection{Where the models fail}
Tables~\ref{tab:errorsA} and~\ref{tab:errorsB} split the test error by
price band, property category and the number of training listings in the
listing's district; the band edges and the district counts come from the
training rows only.

\textit{The luxury tail.} The error on the most expensive 1\% of listings
is several times the error in the middle of the market. A tree of limited
depth has a limited number of leaves, so rare expensive listings share a
leaf with cheaper ones and are pulled towards their mean; the largest
misses in Table~\ref{tab:worst} are dominated by under-predictions of very
expensive land and commercial property.

\textit{Property type.} Flats are numerous and alike, and are predicted
well. Land and commercial property are few, are priced by things the
table does not hold (zoning, frontage, business use), and their area is
on a different scale from that of flats.

\textit{Sparse districts.} Error falls steadily as the number of training
listings in a district grows, in both tasks.

\textit{The tier boundary.} In Task~B the mistakes concentrate in the two
price bands next to the median: a listing a few thousand AZN from the
threshold is close to a coin flip for any model, while the cheapest and
the most expensive listings are classified far more reliably.

\begin{table}[t]
  \centering
  \caption{Task A: mean absolute error on the log price, test split.}
  \label{tab:errorsA}
  \gentable{errors_task_a}
\end{table}

\begin{table}[t]
  \centering
  \caption{Task B: share of wrong tier calls, test split.}
  \label{tab:errorsB}
  \gentable{errors_task_b}
\end{table}

\begin{table}[t]
  \centering
  \caption{The ten largest price errors on the test split, Task A.}
  \label{tab:worst}
  \gentable{worst_misses}
\end{table}

\subsection{Bias and variance}
Fig.~\ref{fig:tuning} shows the usual pattern. Shallow
trees (depth 2--4) score low on both training and validation data: high
bias. As depth grows the training score keeps rising while the validation
score flattens after depth 10, and the widening gap is the tree starting
to memorise its training rows: variance. Raising \texttt{min\_samples\_leaf}
from 1 to 20 lowers the training $R^2$ from 0.95 to 0.80 but keeps the
validation $R^2$ at its best, so larger leaves remove variance at no cost;
at 50 both fall. For the SVM a very small $\lambda$ gives a narrow margin
and a lower validation F1 (0.785 at $10^{-5}$), and a very large one a
wide margin that underfits (0.799 at $10^{-1}$). Note that these sweeps use
a 20{,}000-row sample, so their absolute scores sit below the final ones.

\subsection{Limitations}
\begin{itemize}
  \item Prices are asking prices from listings, not sale prices.
  \item One model serves seven property types. A model per type, or
        dropping land and commercial property, would probably help both
        tasks.
  \item Missing numbers are set to 0. A tree can isolate that value, but a
        linear model reads it as a real floor or land area.
  \item Land plots of up to 50{,}000~m$^2$ give the area a very long tail,
        which hurts the linear SVM after standardising; a logarithm of
        the area would be the first thing to try.
  \item The one-hot columns are created from all listings before the
        split. This uses no target information, but a district seen only
        in the test part still gets a column.
  \item Repeat scrapes are matched by listing URL; a property re-listed
        under a new URL counts as a new listing.
  \item Cross-validation folds share most of their training data, so the
        reported $\pm$ describes stability across folds, not a confidence
        interval.
  \item The data are from early 2025 and almost entirely from Baku.
\end{itemize}

\section{Conclusion}
Both models, written from scratch in NumPy, match
scikit-learn on a real, messy dataset: our tree to the third decimal, our
Pegasos SVM within 0.003 F1, with the gaps explained by tie-breaking,
float precision and the optimiser's stopping rule. A depth-15 tree
predicts Baku asking prices to within about 14\% for ordinary flats and
classifies the price tier with F1 0.91, while a linear SVM is limited by
the joint effect of area and location. Both fail on what is rare: the
luxury tail, land, commercial property and thinly covered districts. Next
we would train a model per property type, use the logarithm of the area
with proper missing-value indicators, and build a random forest from our
own trees.

\section*{Tools and Acknowledgements}
\textit{AI assistance.} F.~Dilavarli designed and led the evaluation side
of the project and coordinated the team's hand-overs. He used Claude Code
(Anthropic) as a coding assistant for the implementation and for drafting
the report and slides, and reviewed its output; the commits concerned carry
a co-author line.
N.~Alimammadova used Claude as an AI assistant to draft the SVM code,
experiments and tests and to explain the method.
N.~Gojayev used Claude (Anthropic) to review the tree, to write the faster
split search, the tests and the comparison script, and to draft report
text; he ran the code and checked the outputs.
U.~Karimova used Claude (Anthropic) as an AI assistant to build the demo
user interface (\texttt{ui/app.py}); she ran it and checked its output on
the real data.

\textit{Data and software.} The dataset is the Bina.az Sale Project
scrape published on Kaggle~\cite{dataset}. We used NumPy, pandas and
matplotlib, and scikit-learn~\cite{sklearn} for the baselines and for
checking our metrics only.

% =====================================================================
\begin{thebibliography}{9}
\bibitem{pegasos}
S.~Shalev-Shwartz, Y.~Singer, N.~Srebro, and A.~Cotter,
``Pegasos: primal estimated sub-gradient solver for SVM,''
\emph{Mathematical Programming}, vol.~127, no.~1, pp.~3--30, 2011.
\bibitem{cart}
L.~Breiman, J.~H. Friedman, R.~A. Olshen, and C.~J. Stone,
\emph{Classification and Regression Trees}. Wadsworth, 1984.
\bibitem{svm}
C.~Cortes and V.~Vapnik, ``Support-vector networks,''
\emph{Machine Learning}, vol.~20, no.~3, pp.~273--297, 1995.
\bibitem{rff}
A.~Rahimi and B.~Recht, ``Random features for large-scale kernel
machines,'' in \emph{Advances in Neural Information Processing Systems
20}, 2007, pp.~1177--1184.
\bibitem{sklearn}
F.~Pedregosa \emph{et al.}, ``Scikit-learn: machine learning in Python,''
\emph{Journal of Machine Learning Research}, vol.~12, pp.~2825--2830, 2011.
\bibitem{dataset}
``Bina.az Sale Project,'' Kaggle dataset published by user
\texttt{sehriyarmemmedli}, file dated Feb.~2025.
[Online]. Available:
\url{https://www.kaggle.com/datasets/sehriyarmemmedli/binaaz-sale-project}
\end{thebibliography}

\end{document}
